% git_hash=b9befbe7, timestamp=2026-07-16 13:32:50 EDT
%%%% Chapter file for Integrating Causality and Probability in ML %%%%
% This chapter file can be compiled standalone or included in the root book.tex

\chapter{Integrating Causality and Probability in ML}
\label{chap:integrating-causality-and-probability-in-ml}

\motto{A prediction says what will happen. A decision says what to do.}

\abstract*{This chapter builds the causal and probabilistic machinery that the previous chapter's business framing left unexplained. It opens by making the prediction-versus-decision distinction formal: a decision target contains an intervention operator, $do(\cdot)$, that a predictive target never needs, and every technique introduced from here on exists to compute that operator from data. The chapter then introduces Judea Pearl's ladder of causation, which separates causal reasoning into three rungs of increasing power: association, intervention, and counterfactual reasoning. Each rung answers a different kind of question and demands different evidence to answer it. The chapter examines the first rung, association, in detail: it is the rung ordinary predictive machine learning occupies, and also the rung where a model can look accurate while remaining blind to confounding and reverse causation. Later chapters climb the remaining rungs, turning the ladder from a taxonomy into a set of computable tools.}

\abstract{A decision target contains an intervention that a predictive target never does. This chapter formalizes that distinction, introduces the three-rung ladder of causation, association, intervention, counterfactual, and examines the first rung in depth, showing why association alone cannot separate a true cause from a confounder.}

% From: '# Bridging Prediction and Decision'
\section{Bridging Prediction and Decision}
\label{chap:bridging-prediction-and-decision}

% From: '## From Costs to Tools'
\subsection{From Costs to Tools}
\label{sec:from-costs-to-tools}

The previous chapter left four costs unaddressed: causality, uncertainty, the business objective, and dynamics. This lesson builds the machinery needed to close the first two of those gaps, starting with the vocabulary that makes a causal claim precise enough to act on.

% From: '* From Costs to Tools: What This Chapter Builds'
\textbf{From Costs to Tools: What This Chapter Builds}

Chapter~\ref{chap:cost-ignoring-causality} and Chapter~\ref{chap:cost-ignoring-uncertainty} catalogued what a purely predictive pipeline misses: a model that only fits $\Pr(Y \mid X)$ cannot say what happens under an intervention, and a point estimate carries no sense of how much to trust it. This lesson does not revisit those failures; it builds the tools that close them, starting with causality and uncertainty, the first two of the four costs surveyed earlier.

The business framing established in that chapter does not change here: \textit{a prediction says what will happen, and a decision says what to do}. What this chapter adds is the machinery sitting underneath that framing, the apparatus that turns a causal claim from a slogan into a quantity that can actually be computed from data. Three pieces of that apparatus organize the rest of this lesson.

\begin{itemize}
  \item \textbf{The ladder of causation}, which formalizes what the informal phrases \emph{``what if''} and \emph{``why''} actually mean, separating association, intervention, and counterfactual reasoning into distinct rungs with distinct evidence requirements.
  \item \textbf{Causal graphs and identification}, which turn a causal question stated in words into a computable quantity, an expression built from the observable variables a dataset actually contains.
  \item \textbf{A probabilistic toolkit}, which turns a single point estimate into a posterior distribution a decision-maker can trust, replacing a bare number with a quantified sense of how confident that number deserves to be.
\end{itemize}

Each piece answers a different half of the causality-and-uncertainty gap: the ladder and the graphs supply the causal half, saying what a claim like \emph{``if we cut price 10\%, revenue rises''} would even mean and what evidence would settle it; the probabilistic toolkit supplies the uncertainty half, attaching a defensible range to whatever the causal machinery concludes. Neither half is optional. A causal estimate reported without uncertainty is the point-estimate failure mode from the previous chapter, dressed up in causal notation; an uncertainty estimate computed over the wrong, purely observational quantity is the correlation-versus-causation failure mode, dressed up in confidence intervals. The rest of this lesson works through each piece of the apparatus in turn, beginning with the ladder of causation.

% From: '* Prediction vs Decision, Formally'
\textbf{Prediction vs Decision, Formally}

Prediction and decision are not the same kind of mathematical object, and treating them as interchangeable is a common source of costly errors, exactly the kind of error the ice-cream-and-drowning and loan-approval cases explored in the previous chapter. Table~\ref{tab:prediction-vs-decision-formally} sets the two side by side across five dimensions.

\begin{table}[!t]
\caption{Prediction and decision compared across target, data, metric, assumption, and how each treats the world.}
\label{tab:prediction-vs-decision-formally}
\begin{tabular}{p{2.0cm}p{4.5cm}p{4.8cm}}
\hline\noalign{\smallskip}
Aspect & Prediction & Decision \\
\noalign{\smallskip}\svhline\noalign{\smallskip}
Target & $\mathbb{E}[Y \mid X]$ & $\arg\max_a \mathbb{E}[U(Y) \mid do(a)]$ \\
Data & Observational only & Observational + experimental \\
Metric & Accuracy, AUC, MSE & Regret, expected utility \\
Assumption & i.i.d., stable world & Models the intervention \\
World & No feedback & Actions change the world \\
\noalign{\smallskip}\hline\noalign{\smallskip}
\end{tabular}
\end{table}

Reading down the columns of Table~\ref{tab:prediction-vs-decision-formally}, a predictive target is a conditional expectation, $\mathbb{E}[Y \mid X]$, estimated from whatever data happened to be logged, evaluated by how closely it tracks the outcome that actually occurred, and reported under the assumption that the data-generating process holds still. A decision target is an optimization over actions, $\arg\max_a \mathbb{E}[U(Y) \mid do(a)]$, and every one of those properties changes with it: it needs data on what happens \emph{under} an action, not merely data on what an action happened to correlate with; it is scored by regret or expected utility rather than by accuracy; and it explicitly models the intervention rather than assuming the world stays fixed while it is measured. The bottom row of the table is the same warning explored at length under dynamics and feedback in Chapter~\ref{chap:cost-ignoring-dynamics-feedback}: a decision, once acted on, changes the very world the next prediction will be scored against, while a prediction, by construction, assumes it will not.

\textbf{Key idea}: a decision target contains a $do(\cdot)$ that a predictive target never does, and essentially everything introduced in this lesson, the ladder of causation, causal graphs, identification, the probabilistic toolkit, exists for one purpose: computing that $do(\cdot)$ from data that was, in the overwhelming majority of real settings, never generated by an experiment in the first place.

% From: '# The Ladder of Causation'
\section{The Ladder of Causation}
\label{chap:the-ladder-of-causation}

% From: '## Three Rungs of Causal Reasoning'
\subsection{Three Rungs of Causal Reasoning}
\label{sec:three-rungs-of-causal-reasoning}

The apparatus promised in the previous section starts with a map of the different kinds of causal question a model might be asked to answer, and of what each kind of question actually demands as evidence.

% From: '* The Ladder of Causation'
\textbf{The Ladder of Causation}

Judea Pearl organizes causal reasoning into three rungs of increasing power, each one strictly more demanding than the last in terms of what it asks a model to know. Table~\ref{tab:ladder-of-causation} lays out the three rungs side by side, with the symbol, the everyday activity, and the question each rung answers.

\begin{table}[!t]
\caption{Pearl's ladder of causation: three rungs of increasing power, from seeing to doing to imagining.}
\label{tab:ladder-of-causation}
\begin{tabular}{p{2.3cm}p{3.4cm}p{2.0cm}p{2.6cm}}
\hline\noalign{\smallskip}
Rung & Symbol & Activity & Question \\
\noalign{\smallskip}\svhline\noalign{\smallskip}
1. Association & $\Pr(Y \mid X)$ & Seeing & What is? \\
2. Intervention & $\Pr(Y \mid do(X), Z)$ & Doing & What if? \\
3. Counterfactual & $\Pr(Y_x \mid x', y')$ & Imagining & Why? \\
\noalign{\smallskip}\hline\noalign{\smallskip}
\end{tabular}
\end{table}

The rungs in Table~\ref{tab:ladder-of-causation} are not three equally valid ways of describing the same reasoning; they are ordered, and climbing from one to the next requires strictly more than a bigger dataset of the same kind. Rung~1, association, asks what seeing $X$ reveals about $Y$, and it is answered entirely from observational data. Rung~2, intervention, asks what happens if $X$ is set by an outside action, $do(X)$, rather than merely observed, holding a set of other variables $Z$ fixed, and answering it generally requires either an experiment or a causal model built from assumptions that no amount of passive observation can verify on its own. Rung~3, counterfactual reasoning, asks the sharpest question of the three: what would have happened to a specific unit, characterized by an observed treatment $x'$ and outcome $y'$, had a different treatment $x$ been assigned instead. This is a strictly harder question than the population-level intervention question, since it requires reasoning about a world that, for that particular unit, never occurred.

\textbf{Key idea}: a model that only fits $\Pr(Y \mid X)$ lives on rung~1, and no amount of additional rung-1 data, however large the dataset, answers a rung-2 or rung-3 question, because the rungs differ in kind, not in degree; more association-level evidence never becomes intervention-level or counterfactual-level evidence on its own.

% From: '* Rung 1: Association'
\textbf{Rung 1: Association}

Rung~1 asks a single, intuitive question: how should seeing $X$ change a belief about $Y$? Formally, this is the conditional distribution $\Pr(Y \mid X)$, and answering it requires nothing more than passive observation of whether $X$ and $Y$ tend to occur together. This is also the entire territory traditional predictive machine learning occupies: a classifier, a regressor, or a forecasting model, however sophisticated its architecture, is fitting some version of $\Pr(Y \mid X)$ and stopping there.

Two everyday examples make the rung concrete. A symptom is informative about a disease in exactly the rung-1 sense: observing the symptom shifts a belief about which disease is present, without anyone needing to know or model the biological mechanism connecting the two. A poll is informative about an election result the same way, since seeing the poll numbers shifts a belief about the likely outcome purely because the two have tended to move together historically, not because the poll causes the result.

The rung has a sharp limitation, and it is the same limitation the previous chapter spent an entire section on: association alone cannot distinguish a true cause from a confounder, and it cannot distinguish a true cause from reverse causation. Seeing that a symptom and a disease move together does not, by itself, say whether the symptom causes complications, whether the disease causes the symptom, or whether some third factor produces both. A model that lives entirely on rung~1 has no way to adjudicate between these explanations, because all of them predict the identical observational pattern; rung-1 evidence is, in this precise sense, silent on the question that rungs~2 and~3 exist to answer. This is precisely the gap Chapter~\ref{chap:cost-ignoring-causality} identified from the business side, confounding and selection bias masquerading as causal effects, and it is why the next rungs of the ladder, not a larger rung-1 dataset, are the tools that close it.
